\documentclass[12pt,a4paper]{article}
\usepackage[utf8]{inputenc}
\usepackage[russian]{babel}
\usepackage{amsmath,amssymb,amsfonts}
\usepackage{geometry}
\geometry{left=2.5cm,right=1.5cm,top=2cm,bottom=2cm}

\title{\textbf{Исследование нейросетевых архитектур над полями Галуа и разработка адаптивного самообучающегося оракула для декодирования кодов БЧХ}}
\author{Журавлев Александр Антонович, гр. 24.Б81-мм \\
Научный руководитель: к. ф.-м. н., доцент Гориховский В. И.}
\date{Санкт-Петербург, 2026}

\begin{document}
\maketitle

\section*{Аннотация}
Настоящая работа посвящена преодолению фундаментального топологического и алгебраического барьера в задаче нейросетевого декодирования циклических кодов Боуза — Чоудхури — Хоквингема (БЧХ). Существующие графовые (GNN) и взвешенные итеративные нейросетевые подходы (Neural Min-Sum) функционируют в непрерывном пространстве вещественных чисел $\mathbb{R}$, что порождает противоречие при попытке аппроксимировать многомерные проверки четности над дискретным полем $\mathbb{F}_2$ на графах Таннера с экстремальной плотностью коротких 4-циклов (5349 циклов для БЧХ(63,45)). Это приводит к затуханию градиентов и эффекту насыщения ошибок (error floor).

В работе исследуется парадигма нейронных сетей над конечными полями Галуа (GF-NN), чьи веса и операции принадлежат полю $\mathbb{F}_{2^m}$. Предложен математический аппарат оптимизации дискретных слоев на базе вероятностных симплексов через быстрое преобразование Уолша--Адамара (FWHT) и суррогатные градиенты (STE). На базе разработанных GF-сетей проектируется масштабируемая архитектура адаптивного самообучающегося оракула с динамическим жизненным циклом: от имитационного обучения на лету до высокоскоростной эксплуатации с гарантированным синдромным fallback к OSD-1. Ожидается достижение границы максимального правдоподобия при радикальном снижении вычислительных затрат.

\section*{Цель и задачи работы}

\textbf{Цель работы:} Проектирование, теоретическое обоснование и программная реализация нейросетевых архитектур над конечными полями Галуа $\mathbb{F}_{2^m}$, а также разработка на их основе адаптивного самообучающегося оракула для декодирования кодов БЧХ с характеристиками на границе максимального правдоподобия.

\subsection*{Задачи на 1-й семестр (Осенний)}
\begin{enumerate}
    \item Разработать математический аппарат дифференцируемых операторов над конечными полями Галуа $\mathbb{F}_{2^m}$, включая вероятностные релаксации на симплексе с использованием быстрого преобразования Уолша--Адамара (FWHT) и методы суррогатных градиентов (Straight-Through Estimator).
    \item Спроектировать архитектуру базовых дискретных слоев (GFDense, GFActivation) с полиномиальными нелинейностями высокой алгебраической иммунности (мультипликативная инверсия $z \mapsto z^{2^m-2}$ и отображения Голда).
    \item Разработать архитектуру нейросетевого декодера синдромов в пространстве полей Галуа, устраняющую проблему затухания градиентов на плотных фактор-графах кода БЧХ(63,45).
    \item Провести сравнительный экспериментальный анализ сходимости и устойчивости разработанной GF-нейросети относительно классических вещественных моделей (MLP, EW-GNN, Neural Min-Sum).
\end{enumerate}

\subsection*{Задачи на 2-й семестр (Весенний)}
\begin{enumerate}
    \setcounter{enumi}{4}
    \item Спроектировать динамическую конечно-автоматную архитектуру самообучающегося оракула (Self-Evolving Oracle) с трехфазным жизненным циклом: холодный старт, эксплуатация нейросети, синдромный самоконтроль с переходом к OSD-1.
    \item Реализовать асинхронный механизм потокового обучения на лету (On-the-fly Self-Supervised Learning) с буферизацией сложных шумовых паттернов.
    \item Провести полномасштабное имитационное моделирование методом Монте-Карло в канале AWGN ($E_b/N_0 \in [2.0, 9.0]$ дБ), оценив показатели BER/FER, задержку и процент срабатывания fallback-механизма.
    \item Выполнить аппаратно-ориентированную оптимизацию и квантование моделей для программно-определяемых радиосистем (SDR) реального времени.
\end{enumerate}

\end{document}
